\section{The boundary index vanishes}\label{sec:vanishing}

We finish by applying the transfer homomorphism to the boundary family
constructed in Section~\ref{sec:persistence}. Only the actual boundary
inclusions $e:B\to D$ and $i:B\to X$ use assembly naturality.

Let $W_j$ be the filled triangulated disc underlying the lattice graph $D_j$,
and put $W=\coprod_{j\ge2}W_j$. Every elementary triangle has vertices at
pairwise distance one, so there is a factorization by continuous proper maps
\[
 P_1(B)=\coprod_j C_j\longrightarrow W\longrightarrow P_1(D).
\]
The second map includes each elementary triangle as a Rips simplex.
At scale one, no simplex of $P_1(D)$ joins different $D_j$.
A compact subset of this topological disjoint union therefore meets only
finitely many of its finite components. Its inverse image in $W$ is a closed
subset of finitely many compact discs, proving properness of the second map;
the first map is proper for the same reason. Each $W_j$ is contractible;
homotopy invariance in the first variable of $\KK$ and
$\KK_1(\C,\C)=K_1(\C)=0$
\cite[Proposition~17.9.1 and Section~19.8.3]{Blackadar} therefore give
\[
 K_1^{\lf}(W)\cong\prod_j\KK_1(C(W_j),\C)
 =\prod_j\KK_1(\C,\C)=0.
\]
Thus the whole family $\beta_1$ dies in $K_1^{\lf}(P_1(D))$, and
\begin{equation}\label{eq:filledzero}
 e_*\beta=0\quad\text{in }KX_1(D).
\end{equation}

Write $\HH_Y=\ell^2(Y)\otimes\HH_0$ for $Y=B,D,X$. The covering isometries
$E:\HH_B\to\HH_D$ and $I:\HH_B\to\HH_X$ of $e$ and $i$ are
\[
 E(\delta_u\otimes h)=\delta_u\otimes h,
 \qquad I(\delta_u\otimes h)=\delta_{(u,1)}\otimes h
 \quad(u\in B).
\]
Put $e_{\mathrm{Roe}}=\Ad E$ and $i_{\mathrm{Roe}}=\Ad I$.
Because the boundary fibers are the
singleton $Q_0$, the averaging isometry of Section~\ref{sec:transfer}
satisfies the exact identity
\[
 VE=I.
\]
It follows on the entire Roe algebra $C^*(B)$ that
\begin{equation}\label{eq:operatorfactor}
 \Phi\circ e_{\mathrm{Roe}}=i_{\mathrm{Roe}}.
\end{equation}

The following diagram displays the naturality square for $e$ and its
operator-algebra continuation. The lower composite equals
$(i_{\mathrm{Roe}})_*$ by \eqref{eq:operatorfactor}:
\[
\begin{tikzpicture}[baseline=(current bounding box.center),>=Stealth]
 \node (a) at (0,1.3) {$KX_1(B)$};
 \node (b) at (3.9,1.3) {$KX_1(D)$};
 \node (c) at (0,0) {$K_1(C^*(B))$};
 \node (d) at (3.9,0) {$K_1(C^*(D))$};
 \node (f) at (7.4,0) {$K_1(C^*(X))$};
 \draw[->] (a)--node[above]{$e_*$}(b);
 \draw[->] (a)--node[left]{$\mu_B$}(c);
 \draw[->] (b)--node[right]{$\mu_D$}(d);
 \draw[->] (c)--node[below]{$(e_{\mathrm{Roe}})_*$}(d);
 \draw[->] (d)--node[below]{$\Phi_*$}(f);
\end{tikzpicture}
\]
Naturality for $i$, followed by \eqref{eq:operatorfactor}, naturality for
$e$, and \eqref{eq:filledzero}, now gives
\begin{align*}
 \mu_X(\alpha)
 &= (i_{\mathrm{Roe}})_*\mu_B(\beta)\\
 &= \Phi_*(e_{\mathrm{Roe}})_*\mu_B(\beta)\\
 &= \Phi_*\mu_D(e_*\beta)=0.
\end{align*}
The class here is exactly $\alpha=i_*\beta$ tested at every Rips scale in
Proposition~\ref{prop:persistence}. That proposition makes it an
infinite-order class, while \eqref{eq:rationalnonzero} makes
$\alpha\otimes1$ nonzero in the rationalized domain. Its rational index
is zero as well. This proves Theorem~\ref{thm:main}.
